% Source: https://artofproblemsolving.com/community/c3h1945887_heronian_triangle
% Stable identity: aops-topic-1945887-q01
Let the other sides be $a\le b\le17$, and put $u=a+b$, $v=b-a$; then $18\le u\le34$, $0\le v$, $u\equiv v\pmod 2$, and $b\le17$ means $u+v\le34$. Heron's formula gives
\[
16K^2=(a+b+17)(a+b-17)(17-a+b)(17+a-b)=(u^2-289)(289-v^2).
\]
\textbf{Parity.} If $u$ were even, then $v$ is even and all four factors $u\pm17$, $17\pm v$ are odd, so $16K^2$ would be odd, impossible. So $u$ and $v$ are odd.

\textbf{Case $u=19$.} Then $u^2-289=72$ and $K^2=\dfrac{9(289-v^2)}{2}$, so $289-v^2$ must be twice a square, with $v$ odd and $v\le15$:
\[
v=1,3,5,\dots,15:\quad 289-v^2=288,280,264,240,208,168,120,64 .
\]
Halving gives $144,140,132,120,104,84,60,32$, and only $144=12^2$ is a square. So $v=1$: $a=9$, $b=10$ and $K^2=9\cdot144$, $K=36$.

\textbf{Case $u\ge21$.} From $v\le34-u$ we get $289-v^2\ge289-(34-u)^2=(u-17)(51-u)$, so
\[
16K^2\ge(u-17)^2(u+17)(51-u).
\]
For $u=21$ this is only $16K^2\ge16\cdot1140$, so look closer: here $K^2=\dfrac{152(289-v^2)}{16}=\dfrac{19(289-v^2)}{2}$ with $v$ odd, $v\le13$; $K<36$ would need $289-v^2<\dfrac{2\cdot1296}{19}<137$, so $v=13$, giving $K^2=19\cdot60=1140$, not a square. For $23\le u\le33$, $(u-17)^2\ge36$ and $(u+17)(51-u)\ge50\cdot18=900$, so $K^2\ge\dfrac{36\cdot900}{16}=2025$ and $K\ge45$.

Hence the least possible area is $\boxed{36}$, attained by the triangle with sides $9,10,17$ (semiperimeter $18$, area $\sqrt{18\cdot9\cdot8\cdot1}=36$).

\begin{asy}
usepackage("amsmath");
size(7cm);
defaultpen(fontsize(10pt));
pair A=(0,0),B=(17,0),C=(308/34,72/17);
draw(A--B--C--cycle,linewidth(1));
draw(C--(C.x,0),dashed);
draw((C.x,0)--(C.x,0.5)--(C.x-0.5,0.5)--(C.x-0.5,0));
label("$A$",A,SW);label("$B$",B,SE);label("$C$",C,N);
label("$17$",(A+B)/2+(-2.5,0),S);label("$10$",(A+C)/2,NW);label("$9$",(B+C)/2,NE);
label("$h=\dfrac{72}{17}$",(C.x,1.6),E);
label("\mbox{}",(min(currentpicture,user=true)+max(currentpicture,user=true))/2,black); // keeps the text after the figure black when LaTeX typesets the labels inline
\end{asy}
